\chapter{Non-collapsed results}\label{chap:noncollapsed}
\label{noncoll}

In this chapter we apply the results for the reduced length function
and reduced volume established in the last two sections to prove
non-collapsing results. In the first section we give a general
result that applies to generalized Ricci flows and will eventually
be applied to Ricci flows with surgery to prove the requisite
non-collapsing. In the second section we give a non-collapsing
result for Ricci flows on compact $3$-manifolds with normalized
initial metrics.

\section{A non-collapsing result for generalized Ricci flows}

The main result of this chapter is a
$\kappa$-non-collapsed result.

 \begin{theorem}[$\kappa$-non-collapsing for generalized Ricci flows]\label{thm:kappa-noncollapsed-generalized-flow}
\label{THM}
 \uses{def:generalized-ricci-flow,def:compatible-embedding}
 Fix  positive constants  $\bar\tau_0<\infty$, $l_0<\infty$, and  $V>0$.
 Then there is $\kappa>0$ depending on $\bar\tau_0$, $V$, and $l_0$
 and the dimension $n$ such
that the following holds. Let $({\mathcal M},G)$ be a generalized
$n$-dimensional Ricci flow, and let $0<\tau_0\le \bar\tau_0$. Let
$x\in {\mathcal M}$ be fixed. Set $T={\bf t}(x)$. Suppose that
$0<r\le \sqrt{\tau_0}$ is given. These data are required to satisfy:
\begin{enumerate} \item[(1)] The ball $B(x,T,r)\subset M_T$ has compact
closure. \item[(2)] There is an embedding $B(x,T,r)\times
[T-r^2,T]\subset {\mathcal M}$ compatible with ${\bf t}$ and with
the vector field. \item[(3)]  $|\Rm|\le r^{-2}$ on the image of
the embedding in (2). \item[(4)] There is an open subset $\widetilde
W\subset \widetilde {\mathcal U}_x(\tau_0)\subset T_x M_T$ with the
property that for every ${\mathcal L}$-geodesic $\gamma\colon
[0,\tau_0]\to {\mathcal M}$ with initial condition contained in
$\widetilde W$, the $l$-length of $\gamma$ is at most $l_0$.
\item[(5)] For each $\tau\in [0,\tau_0]$, let $W(\tau)={\mathcal
L}\operatorname{exp}_x^\tau(\widetilde W)$. The volume of the image
$W(\tau_0)\subset M_{T-\tau_0}$ is at least $V$.
\end{enumerate}
Then
$$\Vol(B(x,T,r))\ge \kappa r^n.$$
 \end{theorem}

In this section we denote by $g(\tau), \ 0\le \tau\le r^2$, the
family of metrics on $B(x,T,r)$ induced from pulling back $G$ under
the embedding $B(x,T,r)\times [T-r^2,T]\to {\mathcal M}$. Of course,
this family of metrics satisfies the backward Ricci flow equation.

\begin{proof}
\uses{def:reduced-volume,thm:reduced-volume-monotone}
Clearly from the definition of the reduced volume, we have
$$\widetilde V_x(W(\tau_0))\ge \tau_0^{-n/2}Ve^{-l_0}\ge
\bar\tau_0^{-n/2}Ve^{-l_0}.$$ By the
monotonicity result (Theorem~\ref{Amono}) it follows that for any
$\tau\le \tau_0$, and in particular for any $\tau\le r^2$, we have
\begin{equation}\label{redvoleqn} \widetilde V_x(W(\tau))\ge
\bar\tau_0^{-n/2}Ve^{-l_0}.\end{equation}

Let $\varepsilon=\sqrt[n]{\Vol(B(x,T,r))}/r$, so that $\Vol\,B(x,T,r)=\varepsilon^nr^n$.  The basic result we need to
establish in order to prove this theorem is the following:


\begin{proposition}[Upper bound for the reduced volume of $W(\tau_1)$]\label{prop:reduced-volume-upper-bound}
\label{step}
\uses{thm:kappa-noncollapsed-generalized-flow}
There is a positive constant
$\varepsilon_0\le 1/4n(n-1)$ depending on $\bar\tau_0$ and $l_0$
such that if $\varepsilon\le \varepsilon_0$ then, setting
$\tau_1=\varepsilon r^2$, we have
 $\widetilde V_x(W(\tau_1)) < 3\varepsilon^{\frac{n}{2}}$.
\end{proposition}

Given this proposition, it follows immediately that either
$\varepsilon>\varepsilon_0$ or
$$\varepsilon\ge \left(\frac{\widetilde V_x(W(\tau_1))}{3}\right)^{2/n}\ge
\frac{1}{3^{2/n}\bar\tau_0}V^{2/n}e^{-2l_0/n}.$$ Since
$\kappa=\varepsilon^n$, this proves the theorem.

\begin{proof}  We divide $\widetilde W$
into
$$\widetilde W_{\mathrm{sm}}=\widetilde W\cap \left\{Z\in
T_xM_T\bigl|\bigr. |Z|\le \frac{1}{8}\varepsilon^{-1/2}\right\}$$
and
$$\widetilde W_{\mathrm{lg}}=\widetilde W\setminus \widetilde W_{\mathrm{sm}},$$

We set $W_{\mathrm{sm}}(\tau_1)={\mathcal L}\operatorname{exp}_x^{\tau_1}(\widetilde W_{\mathrm{sm}})$ and $W_{\mathrm{lg}}(\tau_1)={\mathcal L}\operatorname{exp}_x^{\tau_1}(\widetilde W_{\mathrm{lg}}).$ Clearly, since $W(\tau_1)$ is the union of $W_{\mathrm{sm}}(\tau_1)$ and $W_{\mathrm{lg}}(\tau_1)$ and since these subsets are
disjoint measurable subsets, we have
$$\widetilde V_x(W(\tau_1))=\widetilde V_x(W_{\mathrm{sm}}(\tau_1))+\widetilde V_x(W_{\mathrm{lg}}(\tau_1)).$$  We shall show that there is $\varepsilon_0$ such that either
$\varepsilon>\varepsilon_0$ or  $\widetilde V_x(W_{\mathrm{sm}}(\tau_1))\le
2\varepsilon^{n/2}$ and $\widetilde V_x(W_{\mathrm{lg}}(\tau_1))\le
\varepsilon^{n/2}$. This will establish Proposition~\ref{prop:reduced-volume-upper-bound} and hence
Theorem~\ref{thm:kappa-noncollapsed-generalized-flow}.

\subsection{Upper bound for  $\widetilde V_x(W_{\mathrm{sm}}(\tau_1))$}

The idea here is that ${\mathcal L}$-geodesics with initial vector
in $\widetilde W_{\mathrm{sm}}$ remain in the parabolic neighborhood
$P=B(x,T,r)\times [T-r^2,T]$ for $\tau\in [0,r^2]$. Once we know
this it is easy to see that their ${\mathcal L}$-lengths are bounded
from below. Then if the volume of $B(x,T,r)$ was arbitrarily small,
the reduced volume of  $W_{\mathrm{sm}}(\tau_1)$ would be arbitrarily
small.

\begin{lemma}[Reduced volume bound for the small-vector region]\label{lem:reduced-volume-small-directions}
\label{small}
\uses{prop:reduced-volume-upper-bound}
Setting $\tau_1=\varepsilon r^2$, there is a constant
$\varepsilon_0>0$ depending on $\bar\tau_0$ such that, if
$\varepsilon\le \varepsilon_0$, we have
$$\int_{\widetilde
W_{\mathrm{sm}}}\tau_1^{-n/2}e^{-\tilde l(Z,\tau_1)}{\mathcal
J}(Z,\tau_1)dZ\le 2\varepsilon^{\frac{n}{2}}.$$
\end{lemma}

Of course,  we have
$$\widetilde V_x(W_{\mathrm{sm}}(\tau_1))=\int_{W_{\mathrm{sm}}(\tau_1)}\tau_1^{-n/2}e^{-l(q,\tau_1)}d\vol_{g(\tau_1)}=\int_{\widetilde W_{\mathrm{sm}}}\tau_1^{-n/2}e^{-\tilde
l(Z,\tau_1)}{\mathcal J}(Z,\tau_1)dZ,$$ so that it will follow
immediately from the lemma that:

\begin{corollary}[Reduced volume of $W_{\mathrm{sm}}(\tau_1)$ is small]\label{cor:reduced-volume-small-bound}
\uses{lem:reduced-volume-small-directions}
There is a  constant $\varepsilon_0>0$ depending on $\bar\tau_0$
such that, if $\varepsilon\le \varepsilon_0$, we have
$$\widetilde V_x(W_{\mathrm{sm}}(\tau_1))\le 2\varepsilon^{\frac{n}{2}}.$$
\end{corollary}



\begin{proof}[Proof of Lemma~\ref{lem:reduced-volume-small-directions}]
\uses{claim:l-geodesic-stays-in-ball,claim:volume-comparison-time-slices}
In order to establish Lemma~\ref{lem:reduced-volume-small-directions} we need two preliminary
estimates:


\begin{claim}[Gradient of scalar curvature and metric comparison bounds]\label{claim:scalar-curvature-gradient-metric-bound}
\label{nabR}
\uses{thm:kappa-noncollapsed-generalized-flow}
There is a universal positive constant $\varepsilon'_0$ such that, if
$\varepsilon\le \varepsilon'_0$, then there is a constant $C_1<\infty$
depending only on the dimension $n$ such that the following hold for all $y\in
B(x,T,r/2)$, and for all $t\in[T-\tau_1,T]$:
\begin{enumerate}
\item[(1)] $$|\nabla R(y,t)|\le \frac{C_1}{r^3}$$
\item[(2)]  $$(1-C_1\varepsilon)\le \frac{g(y,t)}{g(y,T)}\le (1+C_1\varepsilon).$$
\end{enumerate}
\end{claim}

\begin{proof}
\uses{thm:shi-first-derivative}
Recall that by hypothesis $|\Rm(y,t)|\le 1/r^2$ on
$B(x,T,r)\times [T-r^2,T]$. Rescale the flow by multiplying the
metric and time by $r^{-2}$ resulting in a ball $\widetilde B$ of
radius one and a flow defined for a time interval of length one with
$|\Rm|\le 1$ on the entire parabolic neighborhood
$B(x,T,1)\times [T-1,T]$. Then according to Theorem~\ref{thm:shi-first-derivative} there
is a universal constant $C_1$ such that $|\nabla R(y,t)|\le C_1$ for
all $(y,t)\in B(x,T,1/2)\times [T-1/2,T]$. Rescaling back by $r^2$
to the original flow, we see that on this flow $|\nabla R(y,t)|\le
C_1/r^3$ for all $(y,t)\in B(x,T,r/2)\times [T-r^2/2,T]$. Taking
$\varepsilon'_0\le 1/2$ gives the first item in the claim.

 Since  $|\Ric|\le (n-1)/r^2$ for all
$(y,t)\in B\times[T-r^2,T]$ it follows by integrating that
$$e^{-2(n-1)(T-t)/r^2}\le \frac{g(x,t)}{g(x,T)}\le e^{2(n-1)(T-t)/r^2}.$$
Thus, for $t\in [T-\tau_1,T]$ we have
$$e^{-2(n-1)\varepsilon}\le \frac{g(x,t)}{g(x,T)}\le e^{2(n-1)\varepsilon}.$$
From this the second item in the claim is immediate.
\end{proof}

At this point we view the ${\mathcal L}$-geodesics as paths
$\gamma\colon [0,\tau_1]\to B(x,T,r)$ (with the understanding that
the path in space-time is given by the composition of the path $
(\gamma(\tau),T-\tau)$ in $B(x,T,r)\times [T-r^2,T]$ followed by the
given inclusion of this product into ${\mathcal M}$.


The next step in the  proof is to show that for any $Z\in \widetilde
W_{\mathrm{sm}}$ the ${\mathcal L}$-geodesic $\gamma_Z$ (the one having
$\lim_{\tau\rightarrow 0}\sqrt\tau X_{\gamma_Z}(\tau)=Z$)
remains in $B(x,T,r/2)$ up to time $\tau_1$. Because of this, as we
shall see, these paths contribute a small amount to the reduced
volume since $B(x,T,r/2)$ has small volume. We set
$X(\tau)=X_{\gamma_Z}(\tau)$


\begin{claim}[$\mathcal L$-geodesic norm comparison estimate]\label{claim:l-geodesic-norm-comparison}
\uses{claim:scalar-curvature-gradient-metric-bound}
There is  a
positive constant $\varepsilon_0\le 1/4n(n-1)$ depending on
$\bar\tau_0$, such that the following holds. Suppose that
 $\varepsilon\le \varepsilon_0$  and $\tau_1' \leq
\tau_1=\varepsilon r^2$. Let $Z\in T_xM_T$ and let $\gamma_Z$ be the
associated ${\mathcal L}$-geodesic from $x$. Suppose  that
$\gamma_Z(\tau) \in B(x,T,r/2)$ for all $\tau < \tau_1'$. Then for
all $\tau<\tau_1'$ we have
$$\abs{|\sqrt{\tau}X(\tau)|_{g(T)} - |Z|} \leq 2\varepsilon (1+|Z|).$$
\end{claim}

\begin{proof}
\uses{claim:scalar-curvature-gradient-metric-bound,lem:l-geodesic-min-max-estimate}
First we make sure that $\varepsilon_0$ is less than or equal to the
universal constant $\varepsilon_0'$ of the last claim. For all
$(y,t)\in B(x,T,r)\times [T-r^2,T]$ we have $|\Rm(y,t)|\le
r^{-2}$ and $|\nabla R(y,t)|\le C_1/r^3$ for some universal constant
$C_1$. Of course, $r^2\le \bar\tau$. Thus, at the expense of
replacing $C_1$ by a larger constant, we can (and shall) assume that
$C_1/r^3>(n-1)r^{-2}\ge |\Ric(y,t)|$ for all $(y,t)\in
B(x,T,r)\times [T-r^2,T]$. Thus, we can take the constant $C_0$ in
the hypothesis of Lemma~\ref{min-max} to be $C_1/r^3$. We take the
constant $\bar\tau$ in the hypothesis of that lemma to be
$\varepsilon r^2$. Then, we have that
$$\max_{0\le \tau\le \tau_1'}\sqrt{\tau}|X(\tau)|\le
e^{2C_1\varepsilon^2}|Z|+\frac{e^{2C_1\varepsilon^2}-1}{2}\sqrt{\varepsilon}
r,$$ and
$$|Z|\le e^{2C_1\varepsilon^2}\min_{0\le \tau\le
\tau_1'}\sqrt{\tau}|X(\tau)|+\frac{e^{2C_1\varepsilon^2}-1}{2}\sqrt{\varepsilon}r.$$
By choosing $\varepsilon_0>0$ sufficiently small (as determined by
the universal constant $C_1$ and by $\bar\tau_0$), we have
$$\max_{0\le \tau\le \tau_1'}\sqrt{\tau}|X(\tau)|_{g(T-\tau)}\le
(1+\frac{\varepsilon}{2})|Z|+\frac{\varepsilon}{2},$$ and
$$|Z|\le (1+\frac{\varepsilon}{2})\min_{0\le \tau\le
\tau_1'}\sqrt{\tau}|X(\tau)|_{g(T-\tau)}+\frac{\varepsilon}{2}.$$ It
is now immediate that
$$\abs{|\sqrt{\tau}X(\tau)|_{g(T-\tau)} - |Z|} \leq \varepsilon (1+|Z|).$$


Again choosing $\varepsilon_0$ sufficiently small the result now
follows from the second inequality in Claim~\ref{claim:scalar-curvature-gradient-metric-bound}
\end{proof}




Now we are ready to establish that the ${\mathcal L}$-geodesics
whose initial conditions are elements of $\widetilde W_{\mathrm{sm}}$ do
not leave $B(x,T,r/2)\times [T-r^2,T]$ for any  $\tau\le\tau_1$.

\begin{claim}[$\mathcal L$-geodesics with small initial vector stay in the half-ball]\label{claim:l-geodesic-stays-in-ball}
\label{inside}
\uses{claim:l-geodesic-norm-comparison}
Suppose $\varepsilon_0\le 1/4n(n-1)$ is the constant from the last
claim. Set $\tau_1=\varepsilon r^2$, and suppose that
$\varepsilon\le \varepsilon_0$. Lastly, assume that $|Z|\le
\frac{1}{8\sqrt{\varepsilon}}$. Then $\gamma_Z(\tau) \in B(x,T,r/2)$
for all $\tau\le \tau_1$.
\end{claim}

\begin{proof}
\uses{claim:l-geodesic-norm-comparison}
Since $\varepsilon\le \varepsilon_0\le 1/4n(n-1)\le 1/8$, by the
last claim we have
$$|\sqrt{\tau}X(\tau)|_{g(T)}\le (1+2\varepsilon )|Z|+2\varepsilon\le
\frac{5}{4}|Z|+\frac{3}{32\sqrt\varepsilon},$$ provided that
$\gamma|_{[0,\tau)}$ is contained in $B(x,T,r/2)\times [T-\tau,T]$.
Since $|Z|\le (8\sqrt\varepsilon)^{-1}$ we conclude that
$$|\sqrt{\tau}X(\tau)|_{g(T)}\le \frac{1}{4\sqrt\varepsilon},$$
as long as $\gamma([0,\tau))$ is contained in $B(x,T,r/2)\times
[T-\tau,T]$.

 Suppose that there is
$\tau'<\tau_1=\varepsilon r^2$ for which $\gamma_Z$ exits
$B(x,T,r/2)\times [T-r^2,T]$. We take $\tau'$ to be the first such
time. Then we have
$$
|\gamma_Z(\tau')- x|_{g(T)} \leq
\int_{0}^{\tau'}|X(\tau)|_{g(T)}d\tau \le
\frac{1}{4\varepsilon^{\frac{1}{2}}}\int_{0}^{\tau'}\frac{d\tau}{\sqrt{\tau}}
=\frac{1}{2\varepsilon^{\frac{1}{2}}}\sqrt{\tau'} < r/2.
$$
This contradiction implies that  $\gamma_Z(\tau) \in B(x,T,r/2)$ for
all $\tau< \tau_1=\varepsilon r^2$.
\end{proof}




Now we assume that $\varepsilon_0>0$ depending on $\bar\tau_0$ is as
above and that $\varepsilon\le \varepsilon_0$, and we shall estimate
$$
\widetilde V_x(W_{\mathrm{sm}}(\tau_1)) = \int_{W_{\mathrm{sm}}(\tau_1)}
(\tau_1)^{-\frac{n}{2}}e^{- l(q,\tau_1)}d\vol_{g(\tau_1)}.$$ In
order to do this we estimate $ l_x(q,\tau_1)$ on $ W_{\mathrm{sm}}(\tau_1)$. By hypothesis $|\Rm|\le 1/r^2$ on
$B(x,T,r/2)\times [0,\tau_1]$ and by Lemma~\ref{claim:l-geodesic-stays-in-ball} every
${\mathcal L}$-geodesic $\gamma_Z$, defined on $[0,\tau_1]$, with
initial conditions $Z$ satisfying $|Z|\leq
\frac{1}{8}\varepsilon^{-\frac{1}{2}}$ remains in $B(x,T,r/2)$.
Thus, for such $\gamma_Z$ we have $R(\gamma_Z(\tau))\ge
-n(n-1)/r^2$. Thus, for any $q\in W_{\mathrm{sm}}(\tau_1)$ we have
$$ L_x(q,\tau_1) =
\int_{0}^{\tau_1}\sqrt{\tau}(R+|X(\tau)|^{2})d \tau \geq
-\frac{2n(n-1)}{3r^{2}}(\tau_1)^{\frac{3}{2}}=
-\frac{2n(n-1)}{3}\varepsilon^{\frac{3}{2}}r,$$ and hence
$$l+x(q,\tau_1)=\frac{L_x(q,\tau_1)}{2\sqrt{\tau_1}}\ge -\frac{n(n-1)}{3}\varepsilon.$$ Since $W_{\mathrm{sm}}(\tau)\subset B(x,T,r/2)\subset B(x,T,r)$, we have:
\begin{eqnarray}\label{volineq}
 \widetilde V_x(W_{\mathrm{sm}}(\tau_1)) & \leq &
\varepsilon^{-\frac{n}{2}}r^{-n}e^{n(n-1)\varepsilon/3}\Vol_{g(T-\tau_1)}\, W_{\mathrm{sm}}(\tau) \\
& \le & \varepsilon^{-\frac{n}{2}}r^{-n}e^{n(n-1)\varepsilon/3}\Vol_{g(T-\tau_1)}\,B(x,T,r).\nonumber\end{eqnarray}








\begin{claim}[Volume comparison between nearby time-slices]\label{claim:volume-comparison-time-slices}
\label{volvomp}
\uses{thm:kappa-noncollapsed-generalized-flow}
There is a universal constant $\varepsilon_0>0$ such that if
$\varepsilon\le \varepsilon_0$, for any open subset $U$ of
$B(x,T,r)$, and for any $0\le \tau_1\le \tau_0$, we have
$$0.9\le \Vol_{g(T)}U/\Vol_{g(T-\tau_1)}U\le 1.1.$$
\end{claim}


\begin{proof}
\uses{claim:scalar-curvature-gradient-metric-bound}
This is immediate from the second item in Claim~\ref{claim:scalar-curvature-gradient-metric-bound}.
\end{proof}

Now assume that $\varepsilon_0$ also satisfies this claim. Plugging
this into Equation~(\ref{volineq}), and using the fact that
$\varepsilon\le \varepsilon_0\le 1/4n(n-1)$, , so that
$n(n-1)\varepsilon/3\le 1/12$, and the fact that from the definition
we have $\Vol_{g(T)}\,B(x,T,r)=\varepsilon^nr^n$, gives
$$\widetilde V_x(W_{\mathrm{sm}}(\tau_1))\leq \varepsilon^{-\frac{n}{2}}r^{-n}e^{n(n-1)\varepsilon/3}
(1.1)\Vol_{g(T)}B(x,T,r)\le(1.1)
\varepsilon^{\frac{n}{2}}e^{\frac{1}{12}}.$$ Thus,
$$\widetilde V_x(W_{\mathrm{sm}}(\tau_1))\le
2\varepsilon^{\frac{n}{2}}.$$

This completes the proof of Lemma~\ref{lem:reduced-volume-small-directions}.
\end{proof}


\subsection{Upper bound for $\widetilde V_x( W_{\mathrm{lg}}(\tau_1))$}

Here the basic point is to approximate the reduced volume integrand
by the heat kernel, which drops off exponentially fast as we go away
from the origin.




Recall that $\Vol\,B(x,T,r)=\varepsilon^nr^n$ and
$\tau_1=\varepsilon r^2$.

\begin{lemma}[Reduced volume bound for the large-vector region]\label{lem:reduced-volume-large-directions}
\label{large}
\uses{prop:reduced-volume-upper-bound}
There is a universal positive constant $\varepsilon_0>0$ such that
if $\varepsilon\le \varepsilon_0$, we have
$$
\widetilde V_x(W_{\mathrm{lg}}(\tau_1)) \le  \int_{{\widetilde
U}(\tau_1)\cap \{Z\bigl|\bigr. |Z|\geq
\frac{1}{8}\varepsilon^{-\frac{1}{2}}\}}(\tau_1)^{-\frac{n}{2}}
e^{-\tilde l(q,\tau_1)}\mathcal{J}(Z,\tau_1)dZ\le
\varepsilon^{\frac{n}{2}}.$$
\end{lemma}

\begin{proof}
\uses{prop:reduced-volume-integrand-monotone}
By the monotonicity result (Proposition~\ref{fmonotone}), we see
that the restriction of the function
$\tau_1^{-\frac{n}{2}}e^{-\tilde l(Z,\tau_1)}{\mathcal J}(Z,\tau_1)$
to ${\widetilde U}(\tau_1)$ is less than or equal to the restriction
of the function $2^ne^{-|Z|^2}$ to the same subset. This means that
$$\widetilde V_x(W_{\mathrm{lg}}(\tau_1))\le \int_{{\widetilde U}(\tau_1)\setminus
\widetilde U(\tau_1)\cap B(0,\frac{1}{8}\varepsilon^{-1/2})}2^n
e^{-|Z|^2}dZ\le \int_{T_p{\mathcal M}_T\setminus
B(0,\frac{1}{8}\varepsilon^{-1/2})}2^ne^{-|Z|^2}dZ.$$ So it suffices
to estimate this latter integral.

Fix some $a>0$ and  let $I(a)=\int_{B(0,a)}2^ne^{-|Z|^2}dZ$. Let
$R(a/\sqrt{n})$ be the $n$-cube centered at the origin with side
lengths $2a/\sqrt{n}$. Then $R(a/\sqrt{n})\subset B(0,a)$, so that
\begin{eqnarray*}
I(a) & \ge  & \int_{R(a/\sqrt{n})}2^ne^{-|Z|^2}dZ \\
& = & \prod_{i=1}^n
\left(\int_{-a/\sqrt{n}}^{a/\sqrt{n}}2e^{-z_i^2}dz_i\right)
\\
& = &
\left(\int_0^{2\pi}\int_0^{a/\sqrt{n}}4e^{-r^2}rdrd\theta\right)^{n/2}.
\end{eqnarray*}

Now
$$\int_0^{2\pi}\int_0^{a/\sqrt{n}}4e^{-r^2}rdrd\theta=4\pi(1-e^{-\frac{a^2}{n}}).$$
Applying this with $a=(8\sqrt\varepsilon)^{-1}$ we have
$$\widetilde V_x(W_{\mathrm{lg}}(\tau_1))\le \int_{\Ar^n}2^ne^{-|Z|^2}dZ-I(1/8\sqrt\varepsilon)
\le(4\pi)^{n/2}\left(1-\left(1-e^{-1/(64n\varepsilon)}\right)^{n/2}\right).$$
Thus,
$$\widetilde V_x(W_{\mathrm{lg}}(\tau_1))\le
(4\pi)^{n/2}\frac{n}{2}e^{-1/(64n\varepsilon)}.$$ There is $\varepsilon_0>0$ such that the expression on the
right-hand side is less than $\varepsilon^{n/2}$ if $\varepsilon\le
\varepsilon_0$. This completes the proof of Lemma~\ref{lem:reduced-volume-large-directions}.
\end{proof}



Putting Lemmas~\ref{lem:reduced-volume-small-directions} and~\ref{lem:reduced-volume-large-directions} together, establishes
Proposition~\ref{prop:reduced-volume-upper-bound}. \end{proof}

As we have already remarked, Proposition~\ref{prop:reduced-volume-upper-bound} immediately
implies Theorem~\ref{thm:kappa-noncollapsed-generalized-flow}. This completes the proof of
Theorem~\ref{thm:kappa-noncollapsed-generalized-flow}.
\end{proof}



\section{Application to compact Ricci flows}

Now let us apply this result to Ricci flows with normalized initial
metrics to show that they are universally $\kappa$-non-collapsed on
any fixed, finite time interval. In this section we specialize to
$3$-dimensional Ricci flows. We do not need this result in what
follows for we shall prove a more delicate result in the context of
Ricci flows with surgery. Still, this result is much simpler and
serves as a paradigm of what will come.

\begin{theorem}[$\kappa$-non-collapsing for compact 3-dimensional Ricci flows]\label{thm:kappa-noncollapsed-compact-3-manifolds}
\uses{def:ricci-flow}
Fix positive constants $\omega>0$ and $T_0<\infty$. Then there is
$\kappa>0$ depending only on these constants such that the following
holds. Let $(M,g(t)),\ 0\le t<T\le T_0$, be a $3$-dimensional Ricci
flow with $M$ compact. Suppose that $|\Rm(p,0)|\le 1$ and also
that $\Vol\,B(p,0,1)\ge \omega$ for all $p\in M$. Then for any
$t_0\le T$, any  $r>0$ with $r^2\le t_0$ and any $(p,t_0)\in M\times
\{t_0\}$, if $|\Rm(q,t)|\le r^{-2}$ on $B(p,t_0,r)\times
[t_0-r^2,t_0]$ then $\Vol\,B(p,t_0,r)\ge \kappa r^3$.
\end{theorem}


\begin{proof}
\uses{thm:bishop-gromov,thm:kappa-noncollapsed-generalized-flow,prop:normalized-flow-short-time-bounds,thm:min-l-bound-n-half,prop:l-exp-lipschitz-complete}
Fix any $x=(p,t_0)\in M\times [0,T]$. First, we claim that we can
suppose that $t_0\ge 1$. For if not, then rescale the flow by
$Q=1/t_0$. This does not affect the curvature inequality at time
zero. Furthermore, there is $\omega'>0$ depending only on $\omega$
such that for any ball $B$ at time zero and of radius one in the
rescaled flow we have $\Vol\,B\ge \omega'$. The reason for the
latter fact is the following: By the Bishop-Gromov inequality
(Theorem~\ref{thm:bishop-gromov}) there is $\omega'>0$ depending only on
$\omega$ such that for any $q\in M$ and any $r\le 1$ we have $
\Vol\,B(q,0,r)\ge \omega'r^3$. Of course, the rescaling increases
$T$, but simply restrict to the rescaled flow on $[0,1]$.

Next, we claim that we can assume that $r\le \sqrt{t_0}/2$. If $r$
does not satisfy this inequality, then we replace $r$ with
$r'=\sqrt{t_0}/2$. Of course, the curvature inequalities hold for
$r'$ if they hold for $r$. Suppose that we have established the
result for $r'$. Then
$$\Vol\,B(p,T,r)\ge \Vol\,B(p,T,r')\ge \kappa(r')^3\ge
\kappa\left(\frac{r}{2}\right)^3=\frac{\kappa}{8}r^3.$$

From now on we assume that $t_0\ge 1$ and $r\le\sqrt{t_0}/2$.
According to Proposition~\ref{kappa0r0t0} for any $(p,t)\in M\times
[0,2^{-4}]$ we have $|\Rm(p,t)|\le 2$ and $
\Vol\,B(p,t,r)\ge \kappa_0 r^3$ for all $r\le 1$.


Once we know that $|\Rm|$ is universally bounded on $M\times
[0,2^{-4}]$ it follows that there is a universal constant $C_1$ such
that $C_1^{-1}g(q,0)\le g(q,t) \le C_1g(q,0)$ for all $q\in M$ and
all $t\in [0,2^{-4}]$. This means that there is a universal constant
$C<\infty $ such that the following holds. For any points $q_0,q\in
M$ with $d_0(q_0,q)\le 1$ let $\gamma_{q_0,q}$ be the path in
$M\times [2^{-5},2^{-4}]$ given by
$$\gamma_{q_0,q}(\tau)=(A_{q_0,q}(\tau),2^{-4}-\tau),\ 0\le \tau\le 2^{-5},$$
where $A_{q_0,q}$ is a shortest $g(0)$-geodesic from $q_0$ to $q$.
Then ${\mathcal L}(\gamma_{q_0,q})\le C$.


By Theorem~\ref{n/2} there is a point $q_0\in M$ and an ${\mathcal
L}$-geodesic $\gamma_0$ from $x=(p,t_0)$ to $(q_0,2^{-4})$ with
$l(\gamma_0)\le 3/2$. Since $t_0\ge 1$, this means that there is a
universal constant $C'<\infty$ such that for each point $q\in
B(q_0,0,1)$ the path which is the composite of $\gamma_0$ followed
by $\gamma_{q_0,q}$ has $\tilde l$-length at most $C'$. Setting
$\tau_0=t_0-2^{-5}$, this implies that $l_x(q,\tau_0)\le C'$ for
every $q\in B(q_0,0,1)$. This ball has volume at least $\kappa_0$.
By Proposition~\ref{lipcomplete}, the open subset ${\mathcal
U}_x(\tau_0)$ is of full measure in $M\times\{2^{-5}\}$. Hence,
$W(\tau_0)=\left(B(q_0,0,1)\times \{2^{-5}\}\right)\cap {\mathcal
U}_x(\tau_0)$ also has volume at least $\kappa_0$. Since $r^2\le
t_0/4<\tau_0$, Theorem~\ref{thm:kappa-noncollapsed-generalized-flow} now gives the result.
\end{proof}
